% Budget: 10 pp.  Function: what the system IS.  No results, no motivation.
% Carries D1 + the curvature derivation (D3 theory).
\chapter{Method}
\label{ch:method}

\section{The composition interface}
\label{sec:composition}

Everything in this thesis is an instance of one rule. A frozen pretrained
model $\fbase$ is composed with a trainable adapter $\adapter$ through a
gate $\gate$:
\begin{equation}
  \label{eq:composition}
  \composition .
\end{equation}

The base is evaluated with its own inputs only, the state $\xt$ and the
diffusion time $t$, and its parameters are never updated. The adapter
receives those same inputs together with the two conditioning variables the
base does not have: the action $\at$ and the sampling step size $\dstep$.
Its output is a residual on the base's prediction rather than a replacement
for it, and the gate scales that residual as a function of the step size.

Three properties of \cref{eq:composition} are worth drawing out, because the
rest of the thesis turns on them.

\paragraph{The base is a black box.} Only $\fbase(\xt,t)$ is required, not
its weights, its activations or its gradients. This is what makes the method
applicable to models that are released for sampling but not for training,
and it is the constraint under which \cref{sec:families} selects an adapter
family.

\paragraph{The composition is additive and gated, not a mixture.} An
alternative is to combine the two predictions convexly, as
\citet{rigter2024avid} do with a learned mask $m$, giving
$\hat{y} = \hat{y}_{\text{base}} \odot m + \hat{y}_{\text{adapt}} \odot
(1-m)$. We use an additive residual with a multiplicative gate instead. The
difference matters for a specific reason: $\gate$ is a function of
$\dstep$, so the amount by which the adapter is allowed to move the
prediction can itself depend on the size of the step being taken. A convex
mixture of two predictions does not naturally accommodate that, and
step-size conditioning is what \cref{sec:shortcut} needs to express.

\paragraph{The adapter's conditioning surface is closed.} The variables
$(\at, \dstep)$ are what $\adapter$ is given. The base may consume further
inputs of its own; those are discussed in \cref{sec:conditioning} and they
lie outside this interface by construction.

Sections~\ref{sec:conditioning} and~\ref{sec:shortcut} specify the two
conditioning inputs, and \cref{sec:cost} reports what
\cref{eq:composition} costs to train.

\section{Adapter families and the selection}
\label{sec:families}

The framework implements four families of adapter behind
\cref{eq:composition}: low-rank weight updates, hidden-state adapters,
hypernetworks that emit weights, and output-level residual adapters. All of
the empirical work in this thesis uses the last. This section derives that
choice, because a taxonomy that is implemented and then not used is
decoration, and because the criterion that selects it is the same one that
explains a cost result in \cref{sec:res-d1}.

\subsection{Two criteria}

\paragraph{Does the family require access to the base's internals?}
Low-rank updates modify the base's weight matrices; hidden-state adapters
read and write intermediate activations; hypernetworks emit parameters for
some part of the base. Each therefore requires the base to be open in a way
that goes beyond being callable. An output-level adapter requires only
$\fbase(\xt, t)$, and is the only family that survives the black-box
constraint of \cref{sec:composition}. This is
\citeauthor{rigter2024avid}'s argument as well, and it is a claim about
the direction the field is moving in rather than about any particular model:
the strongest video generators are increasingly released as weights that can
be sampled but not inspected.

\paragraph{Does training require gradients through the frozen base?} This
criterion is less obvious and it is the one with a measurable cost. An
output-level adapter can be trained with the base's prediction detached: the
base is evaluated once under \code{no\_grad}, and the backward pass touches
only the adapter. The interior families cannot. Because their trainable
parameters sit inside the base, the forward pass must be differentiable
through it, and the activation memory that this retains then forces
gradient checkpointing at video scale.

The consequence is measured rather than argued. At matched effective batch
on the same frozen base, a low-rank adapter costs \num{3.1} times the
wall-clock per optimiser step of an output-level adapter, running at
\num{66} against \num{205} steps per hour, while being the \emph{smaller}
module of the two at \SI{26.0}{\mega\nothing} against
\SI{34.9}{\mega\nothing} trainable parameters.\prov{25183721}{\nv{ckpt}}{\nv{commit}}
Parameter count is the wrong axis: where the parameters live dominates what
they cost.

\subsection{Why four families are implemented}

If the output adapter wins on both criteria, the natural question is why the
other three exist in the framework at all. Two reasons, and they are
different in kind.

The first is that a selection is only a result if the alternatives were
available. Implementing the taxonomy behind a single interface is what turns
``we used an output adapter'' into ``an output adapter is what the
constraints select'', and it is what allowed the cost measurement above to
be made on a matched setup rather than quoted from another paper.

The second is that the interface is the contribution of
\cref{sec:res-d1}, and an interface with one implementation demonstrates
nothing. The claim that swapping the frozen base is a configuration change
is only meaningful if swapping the adapter is one too.

\subsection{A tension worth stating}

There is an argument against our own choice and it should be made here
rather than left for a reader to find. A weight update that depends on the
step size is, definitionally, a hypernetwork: a network emitting parameters
as a function of $\dstep$. So the hypernetwork family is arguably the
natural home for \cref{sec:shortcut}, and by selecting the output adapter we
are conditioning on $\dstep$ as an \emph{input} rather than realising it as
a \emph{weight update}.

We make that choice on the cost grounds above rather than on expressiveness
grounds, and we do not claim the two are equivalent. What can be said is
that the second criterion applies to a $\dstep$-dependent hypernetwork with
full force: it must emit parameters for something inside the base, so
training it requires differentiating through the base at every step, at the
scale where \cref{sec:cost} shows the base already dominates. That is a
statement about tractability, not about which formulation is better, and we
mark it as a limitation in \cref{sec:limitations}.

\section{Conditioning interfaces}
\label{sec:conditioning}

\Cref{eq:composition} lists the framework's conditioning variables
explicitly: the state $\xt$, the time $t$, the action $\at$ and the step
size $\dstep$. The action enters the adapter as an embedded vector; the step
size enters both the adapter and the gate, and is specified in
\cref{sec:shortcut}. What requires comment is a third input that the base
consumes and the adapter does not.

\subsection{The frame-stride boundary}

A frozen base may take inputs that are part of \emph{its own} training
contract rather than of ours. In the DynamiCrafter case there is a
frame-stride scalar, routed into the base U-Net's time embedding through a
learned embedding network. These base-private inputs sit \textbf{outside}
the framework's conditioning surface: they reach $\fbase$ and do not enter
$\adapter$.

For the adapter wirings we audited, the frame stride reaches the frozen base
only. The adapter's explicit conditioning inputs remain $(\at, \dstep)$; any
temporal-resolution information $\adapter$ consumes is \emph{indirect},
mediated by the base activations it reads through its feature hooks.

We adopt this as a deliberate framework choice rather than an oversight. The
composition rule is action- and step-conditioned, not frame-rate
conditioned. Were a later result to motivate making the frame stride an
explicit input to $\adapter$, that would extend \cref{eq:composition} and is
better treated as a change to the interface than as a configuration detail.

One consequence is worth recording because it affects how the datasets are
described: the numerical value fed to the base's frame-stride channel is a
data-side property, not a framework parameter, and it is specified in
\cref{sec:datasets}.

\section{Shortcut training}
\label{sec:shortcut}

Step-size conditioning is trained with a self-consistency objective: the
model's prediction for a step of size $2\dstep$ is supervised against a
target built by composing two steps of size $\dstep$. This section specifies
how that target is constructed, because the construction is where
\cref{sec:curvature} finds a problem.

Both constructions below are written for diffusion with a velocity
parameterisation (\cfg{model\_type} \code{diffusion},
\cfg{prediction\_type} \code{velocity}) and both rely on a single-step DDIM
micro-step under the training schedule. Per-sample timesteps are supported,
so a batch may carry heterogeneous $(t, t')$ pairs without scheduler-side
bookkeeping.

\paragraph{Base-anchored two-step.} Given the frozen base velocity
$\vel_0 = \fbase(\xt, t)$, advance one schedule step to
$t' = \max(0,\, t-1)$ by a DDIM micro-step, evaluate the base again at the
advanced state and time, $\vel_1 = \fbase(x_{\text{mid}}, t')$, and form
\begin{equation}
  \label{eq:base-anchored}
  y_{\text{target}} = \tfrac{1}{2}\,(\vel_0 + \vel_1).
\end{equation}
This is the average of velocities at the two endpoints, anchored on the
frozen base. It carries no risk of collapse, because the target does not
depend on the adapter; correspondingly, the adapter's step-size input is
decorative under this mode, since the supervision does not vary with it. The
construction uses the state and time pair the adapter sees at the training
step, so target and prediction live in the same schedule frame.

\paragraph{Self-consistency.} Given a sampled step size $\dstep$, the target
is the average of two calls of the \emph{adapted} model at the half step,
chained across one micro-step of that size:
\begin{align}
  \vel_1 &= \fadapted(\xt,\, t,\, \dstep/2), \nonumber \\
  x_{\text{mid}} &= \mathrm{ddim}(\xt,\, \vel_1,\, t,\, t-\dstep), \nonumber \\
  \vel_2 &= \fadapted(x_{\text{mid}},\, t-\dstep,\, \dstep/2), \nonumber \\
  y_{\text{target}} &= \tfrac{1}{2}\,(\vel_1 + \vel_2),
  \label{eq:self-consistency}
\end{align}
with both model calls taken without gradient. Here the adapter is its own
teacher and the frozen base contributes only implicitly, through the
composition inside $\fadapted$. Because a self-supervised target of this
form admits degenerate solutions, anchoring at the finest step is provided
by a standard diffusion loss applied with a fixed probability.

\paragraph{Both constructions honour the temporal advance.} The second
velocity is sampled at the \emph{advanced} time, not at $t$. This is the
single mathematical statement the self-distillation relies on, and an
earlier intermediate implementation of ours evaluated the second call at the
same $t$; we record it here so that the claim to evaluate against the
published target is unambiguous.

\Cref{eq:self-consistency} is the published formulation, and
\cref{sec:curvature} shows that it is exact only under a condition our first
base does not satisfy.

\section{The curvature bias in the shortcut target}
\label{sec:curvature}

The self-consistency objective of \cref{sec:shortcut} supervises a
double-step prediction with the composition of two single steps. Written for
a velocity parameterisation, the target for step size $2\dstep$ is the
average of the velocity at the current state and the velocity at the state
reached after one step of size $\dstep$:
\begin{equation}
  \label{eq:shortcut-target}
  \vel_{\text{target}}(\xt, t, 2\dstep)
  \;=\; \tfrac{1}{2}\Big[\, \vel(\xt, t, \dstep) \;+\; \vel(x_{t+\dstep}, t+\dstep, \dstep) \,\Big].
\end{equation}

This section establishes when \cref{eq:shortcut-target} is the correct
target and when it is not. The distinction is a property of the base
model's probability-flow trajectory, and it determines which bases the
shortcut objective can be applied to as published.

\subsection{Exactness on a straight interpolant}

Under a rectified-flow construction the interpolant between noise and data
is a straight line, and the velocity field along a trajectory is constant.
Two half-steps therefore carry the same velocity as one double-step, and
\cref{eq:shortcut-target} is exact: averaging two identical quantities
returns the quantity. The trajectory has curvature $\curv = 0$, and the
average of two points on a straight line lies on that line.

\subsection{The bias on a curved trajectory}

A variance-preserving diffusion process does not have this property. Its
probability-flow trajectory is curved, and the two velocities being averaged
are tangents at two distinct points of an arc. Their arithmetic mean is the
\emph{Euclidean} centroid of the two tangent directions, whereas the
quantity the objective requires is the direction of the chord that actually
connects the endpoints, a Riemannian mean along the arc.

The discrepancy between the two is the \emph{sagitta} of the arc: for a
trajectory of curvature $\curv$ traversed over a step $\delta$, the chord
departs from the averaged tangent by
\begin{equation}
  \label{eq:sagitta}
  \text{sagitta} \;\approx\; \tfrac{1}{2}\,\curv\,\delta^{2},
\end{equation}
so the error is second order in the step size and vanishes only as
$\curv \to 0$. \Cref{eq:shortcut-target} is therefore not an approximation
that improves with training; it is a systematically wrong target whenever
$\curv \neq 0$, and it is wrong by an amount that grows quadratically with
the step being taken.

\subsection{Why the bias does not train away}

The consequence that matters is not the size of the error at one step but
its behaviour under the doubling tower the objective constructs. A training
objective whose target is biased can still converge to the correct field if
the true field is a \emph{fixed point} of the target-construction rule:
that is, if substituting the true velocities into \cref{eq:shortcut-target}
returns the true double-step velocity. On a curved trajectory it does not:
substituting the exact tangents returns the averaged tangent, which differs
from the chord by \cref{eq:sagitta}.

The true field is therefore not a fixed point of the rule. Each rung of the
doubling ladder is supervised against a target constructed from the rung
below, so the discrepancy is re-injected at every level rather than being
corrected, and it compounds upward. Training longer does not remove it,
because the objective's optimum is not the true field.

\subsection{Numerical verification at zero model error}

The argument above is about the target construction, not about any model, so
it can be checked with the model removed. Substituting exact velocities into
the two candidate constructions under the training schedule, so that the
only error present is the construction's own, the velocity-averaging
target of \cref{eq:shortcut-target} departs from the true displacement by
\SI{5.1}{\percent}, \SI{16.1}{\percent} and \SI{24.1}{\percent} at
$s = \nicefrac{1}{4}$, $\nicefrac{1}{2}$ and $\nicefrac{3}{4}$ respectively,
growing with the step size as \cref{eq:sagitta} predicts. The
endpoint-inversion construction returns $0.000000$ at every step.

\subsection{Two escapes, and the precise form of the claim}

The bias has two independent remedies, corresponding to the two terms in
\cref{eq:sagitta}.

\emph{Fix the coordinates.} Constructing the target by inverting the
endpoints rather than averaging tangents removes the discrepancy without
changing the base, since the chord is then computed directly. This is exact
on the curved manifold, as the numerical check confirms.

\emph{Fix the geometry.} Adopting a base whose trajectory is straight sets
$\curv = 0$ and makes \cref{eq:shortcut-target} faithful as published. This
is the property that motivates the move to a flow-matching base in
\cref{sec:backbones}.

\medskip

We state the claim precisely, because a looser version is both false and
easy to refute. The result is \emph{not} that shortcut modelling fails on
diffusion: consistency distillation on variance-preserving diffusion
demonstrably works, and the endpoint construction above is exact on the same
curved manifold. The result is that \textbf{the published
velocity-averaging target is exact only at zero curvature, and carries a
second-order bias that compounds up the doubling tower otherwise.} That is a
statement about a target construction, and it is what the derivation and the
numerical check jointly support.

\section{Computational profile of the composition}
\label{sec:cost}

\Cref{eq:composition} has a direct consequence for training cost, and it
shapes two framework-level decisions: the frozen base is evaluated on
\emph{every} optimiser step, and the pixel-to-latent encode is factored
\emph{out} of the training loop. Both are justified below.

The measurements are from the flow-based base with the output adapter on one
environment, at \SI{768}{\pixel} square with 65-frame windows and batch size
one, profiled with CUDA-synchronised per-phase timers. They should be read
as characterising this operating point, not as
architecture-independent constants.

\subsection{The frozen base dominates the step}

Because $\adapter$ is a residual \emph{on} $\fbase$, the base forward cannot
be cached across steps: it depends on the freshly sampled $(\xt, t)$, so it
is recomputed every step under \code{no\_grad}. Measured per step on the
cached-latent path: base forward \SI{480}{\milli\second}, adapter forward
\SI{28}{\milli\second}, backward \SI{54}{\milli\second}, optimiser
\SI{3}{\milli\second}, for a total step of
\SI{580}{\milli\second}.\prov{8cug8wfq}{\nv{ckpt}}{\nv{commit}}

The frozen base is therefore about \SI{83}{\percent} of the step, while the
trainable adapter, the only part carrying gradients, is under a fifth. This
is the defining cost signature of adapting a large frozen prior:
\textbf{parameter efficiency does not translate into a proportionally cheap
step}, because the frozen forward is on the critical path. Here the adapter
is \SI{34.97}{\mega\nothing} trainable of \SI{5.03}{\giga\nothing} total,
\SI{0.69}{\percent} of parameters and roughly \SI{17}{\percent} of the time.

The practical implication, taken up in \cref{sec:limitations}, is that the
throughput ceiling lives on the base side: any speedup must target the
frozen forward through kernels or compilation, not the adapter.

\subsection{Latent precompute is a throughput decision, not a memory one}

At this geometry the online pixel-to-latent encode costs approximately
\SI{3.66}{\second} per clip, against the \SI{0.58}{\second} training
step.\prov{hswppa8s}{\nv{ckpt}}{\nv{commit}} The encode alone is thus about
six times the entire step, and online encoding runs at roughly
\num{0.24} steps per second against \num{1.44} from a warm latent cache.

The encode transient coexists comfortably in memory with the resident base,
peaking at about \SI{26}{\giga\byte} reserved on an \SI{80}{\giga\byte}
device, so unlike an earlier native-resolution regime in which precompute
was an out-of-memory workaround, the justification here is purely
wall-clock. We therefore keep an offline latent-precompute stage as the
framework default and treat online encoding as a fallback for cases a static
cache cannot cover, such as unbounded random-window sampling.
